\section{Discussion}
\label{sec:discussion}

\subsection{Implications}

\namei treats VFS name resolution as a policy boundary rather than as a new
filesystem. The evaluation tests whether this boundary lets policy move into a
narrow, verifier-bounded hook without forcing the
developer to own a daemon, a namespace materializer, or a filesystem
implementation.

\subsection{Scope Of The Hypothesis}

\namei targets workloads whose oracle-relevant behavior is path-view policy,
meaning which existing object a name denotes, or whether an object is visible,
while ordinary file semantics stay with the lower filesystem. The evaluation uses
representative workloads that exercise this hypothesis directly and treats
non-path-view source behavior only as workload-boundary context.

\subsection{Extension Path}

The same boundary also defines future extensions. Final-file target selection
and synthetic aliases require additional kernel support, while synthetic
contents, custom metadata persistence, distributed indexing, and cross-path
transactions remain filesystem or runtime responsibilities. \namei is intended
only for workloads whose correctness oracle is path selection or visibility over
existing lower objects.
